Before the formal discussion of this question. Let's consider a 
question in daily life: why the McDonald's and KFC are always together
on a street. To solve this problem, we should establish a mathematical
model. Assume that the population density on this street is invariant
and everyone would choose the closest restaurant to have meal ( If the
two restaurants are at the same point, people will choose one
randomly). If you are the owner of McDonald's and you want to build
a restaurant first before KFC. You want to get as more customers
as you can, which place will you choose?
	
As the picture shown behind, if Mc and KFC open their restaurant at $A$
and $B$. Let $C$ be the middle point of the segment $AB$. We can easily
prove that the customers in one side of C would go to Mc and others
would go to KFC. So if you want more customers than your competitor,
you should choose the middle point of the street or he could open his
restaurant at the middle point of the street and get more than half customers.
	
\begin{center}
    \begin{tikzpicture}
	\draw[ultra thick][black](-3 , 0)--(0,0);
	\draw[ultra thick][gray](0,0)--(3,0);
	\draw [fill] [gray](0,0) circle [radius=0.05];
	\draw (0,0) circle [radius=0.05];
	\node at (0,-0.35) [thick]{$C$};
	\draw [fill] [gray](-1,0) circle [radius=0.05];
	\draw (-1,0) circle [radius=0.05];
	\node at (-1,-0.35) [thick]{$A$};
	\draw [fill] [gray](1,0) circle [radius=0.05];
	\node at (1,-0.35) [thick]{$B$};
	\draw (1,0) circle [radius=0.05];
	\end{tikzpicture}
\end{center}
	
If you want to open $n$ restaurants on the street, it can be proved that
the best choice is dividing the street into $n$ equal parts and choose
middle point of each part.
	
\begin{center}
	\begin{tikzpicture}[scale=2]
		\draw[thick][black](-3 , 0)--(-2,0);
		\draw[thick][gray](-2,0)--(-1.25,0);
		\draw[thick][black](-1.25 , 0)--(-0.55,0);
		\draw[thick][gray](-0.55,0)--(0.5,0);
		\draw[thick][black](0.5 , 0)--(2.35,0);
		\draw[thick][gray](2.35,0)--(3,0);
		
		\draw[fill][gray](-3,0) circle [radius=0.018];
		\draw[fill][gray](-2,0) circle [radius=0.018];
		\draw[fill][gray](-1.25,0) circle [radius=0.018];
		\draw[fill][gray](-0.55,0) circle [radius=0.018];
		\draw[fill][gray](0.5,0) circle [radius=0.018];
		\draw[fill][gray](2.35,0) circle [radius=0.018];
		\draw[fill][gray](3,0) circle [radius=0.018];
		
		\draw [fill] [lightgray](-2.25,0) circle [radius=0.03];
		\node at (-2.25,-0.15) {$A_1$};
		\draw [fill] [lightgray](-0.75,0) circle [radius=0.03];
		\node at (-0.75,-0.15) {$A_2$};
		\draw [fill] [lightgray](0.75,0) circle [radius=0.03];
		\node at (0.75,-0.15) {$A_3$};
		\draw [fill] [lightgray](2.25,0) circle [radius=0.03];
		\node at (2.25,-0.15) {$A_4$};
		
		\draw [fill] [lightgray](-1.75,0) circle [radius=0.03];
		\node at (-1.75,-0.15) {$B_1$};
		\draw [fill] [lightgray](-0.35,0) circle [radius=0.03];
		\node at (-0.35,-0.15) {$B_2$};
		\draw [fill] [lightgray](0.25,0) circle [radius=0.03];
		\node at (0.25,-0.15){$B_3$};
		\draw [fill] [lightgray](2.45,0) circle [radius=0.03];
		\node at (2.45,-0.15) {$B_4$};
		
		\draw (-2.25,0) circle [radius=0.03];
		\draw (-0.75,0) circle [radius=0.03];
		\draw (0.75,0) circle [radius=0.03];
		\draw (2.25,0) circle [radius=0.03];
		\draw (-1.75,0) circle [radius=0.03];
		\draw (-0.35,0) circle [radius=0.03];
		\draw (0.25,0) circle [radius=0.03];
		\draw (2.45,0) circle [radius=0.03];
	\end{tikzpicture}
\end{center}
	
So we naturally want to know that what will happen if we change the
street into the whole city. Could the first one always make sure
he get at least half of the customers?
	
We express this question in another way:
	
For a convex set $T$ in $\R^2$ and a constant $n$, can we find a set
of points $A = \{ A_1, A_2 ... A_n \}$ of $n$ points in $T$   which
satisfy the following conditions?
For any set $B = \{ B_1, B_2 ... B_n \} $ of $n$ points on the plane
, we color $T$ by two colors. For a point $s$ in $T$, find the closet
point to it in $A\cup B$. If it is a point in $A$, color it in red,
otherwise color it in green. The area of red parts in the diagram
is always not less than the green parts. (Here the metric is always
the Euclidean metric)
	
In this article we will try to find out all pairs of $(T, n)$ which meets
the requirement. At first we would give some lemmas in the second
section. Then we would consider a special case ($T$ is a square)
as an example in the third section. In the forth section, we will
give out two strong conclusion to solve the case when $n$ is big enough.
		
We always use capital letters for points and domains, lower-case lette
rs for lines,  $|AB| := d(A, B)$. For a domain $T$ we always use $S_T$
to denote the area of $T$.